\begin{abstract}
En este proyecto se diseñó un sistema fotovoltaico aislado que respalda las cargas críticas de una vivienda en Santa Rosa de Oreamuno, Cartago, servida por JASEC. Las cargas críticas son el refrigerador, el congelador, la iluminación, los cargadores de celular y de laptop y el módem, que consumen \SI{4.33}{kWh} por día. Con el procedimiento del curso, dos días de autonomía, un banco de \SI{24}{V} y una profundidad de descarga de 70\,\%, se requieren \SI{758}{Ah}, que se cubren con cuatro baterías LiFePO\textsubscript{4} Power Queen de \SI{24}{V} y \SI{200}{Ah} en paralelo. El generador tiene cuatro módulos AIKO de \SI{450}{W} (\SI{1.8}{kWp}) en dos cadenas de dos, sobre la cara este del techo, inclinada \SI{12}{\degree}, con un regulador MPPT Victron SmartSolar 150/85 y un inversor-cargador Victron MultiPlus-II de \SI{3000}{VA} con transferencia automática. Incluso en noviembre, el mes de menor radiación, la energía disponible cubre el 116\,\% de la demanda. El diseño incluye la estructura, el cableado, las protecciones y el diagrama unifilar. El precio estimado es de ¢7.10 millones; el sistema reduce la factura anual en ¢128\,165 (54\,\%), pero con los reemplazos de equipos y baterías la inversión no se recupera en 25 años, por lo que su justificación es la continuidad del servicio y no el ahorro.
\end{abstract}

\begin{IEEEkeywords}
sistema fotovoltaico aislado, respaldo, baterías LiFePO\textsubscript{4}, regulador MPPT, inversor-cargador
\end{IEEEkeywords}
